\documentclass[10pt,a4paper]{amsart}
\usepackage[margin=19mm]{geometry}
\usepackage{amsmath,amssymb,mathtools}
\usepackage[scr=rsfs,cal=euler]{mathalfa}
\usepackage{xfrac}
\usepackage[unicode,hidelinks,bookmarks=false]{hyperref}
\newcommand{\ie}{i.e.\ }
\newcommand{\cf}{cf.\ }
\newcommand{\resp}{respectively\ }
\newcommand{\Subsection}{\subsection}
\providecommand{\nobreakdash}{\nobreak\hbox{-}}
\setlength{\emergencystretch}{2em}
\multlinegap0pt
\usepackage{bm}
\usepackage{bbm}
\usepackage{dsfont}
\usepackage{xspace}
\usepackage{cancel}
\usepackage{mathtools}
\usepackage{amsthm}
\makeatletter
\@ifundefined{thm}{\newtheorem{thm}{Theorem}}{}
\@ifundefined{lem}{\newtheorem{lem}[thm]{Lemma}}{}
\@ifundefined{prop}{\newtheorem{prop}[thm]{Proposition}}{}
\makeatother
\title[Finite-dimensional irreducible representations of twisted lo]{Finite-dimensional irreducible representations of twisted loop algebras of the second kind}
\author{Hideya Watanabe}
\date{Source: arXiv:2506.02379v1 (CC BY 4.0)}
\begin{document}
\maketitle

\noindent\emph{Source and licence.} Finite-dimensional irreducible representations of twisted loop algebras of the second kind. Version arXiv:2506.02379v1 (CC BY 4.0), distributed under the Creative Commons Attribution 4.0 International licence, as recorded in the arXiv licence metadata for this exact version and shown on the abstract page https://arxiv.org/abs/2506.02379v1. Full rights notice: licenses/arxiv-2506.02379-CC-BY-4.0.txt.\par\smallskip

\noindent\textit{Editorial scope.} Counted unit: the Thm whose source label is \texttt{thm: classification AI2} in the article (source0.tex lines 1739--1749, source label \texttt{thm: classification AI2}), delivered together with the article's own complete original proof of it (source0.tex lines 1750--1887, one \texttt{proof} environment of 138 lines). No mathematics is written, repaired or re-derived: statement and proof are the article's own text line for line. Required context retained uncounted: prop: coincidence of p' (lines 633--651); thm: sufficient condition for evaluation2 (lines 1081--1090); thm: classification AI1 (lines 1439--1448); lem: deduction from A1 (lines 1702--1712); lem: deduction from AI1 (lines 1718--1732). Every definition, display and label of those lines is retained unchanged. Omitted from this package and not redistributed: 1--632, 652--1080, 1091--1438, 1449--1701, 1713--1717, 1733--1738, 1888--2767. Nothing inside a retained excerpt is omitted or altered: every retained body is delivered exactly as the article's lines are written, so no omission receipt of any kind accompanies this package. Citations: the retained excerpts carry no citation command at all, so no bibliography entry of the article is transcribed and no reference is redistributed. Reference adaptation: none was needed and none was made; every reference inside the delivered unit resolves inside it, so no unresolved reference or citation is produced. Printed slips kept literal: no printed word, symbol, hypothesis or number of the counted statement or of its proof was altered. Layout and class: the article's own class is \texttt{amsart} with packages including amsmath, amsthm, amssymb, color, geometry, tikz, hyperref; the wrapper is the lane's pinned \texttt{amsart} editorial wrapper at 10pt on A4, which restates the theorem-like environments the retained text needs and adds the article's own macro definitions that the retained text needs (none), with extra wrapper packages (bm, bbm, dsfont, xspace, cancel, mathtools). The delivered numbering is the unit's own. Assets: no retained span contains a figure environment, a graphics-inclusion command, a PDF-page inclusion, a table or a TikZ picture; no image file is copied into this package, the package declares no asset dependency and no image file is redistributed with it. Licence: version arXiv:2506.02379v1 of this article is distributed under the Creative Commons Attribution 4.0 International licence, as recorded in the arXiv licence metadata for this exact version and shown on the abstract page https://arxiv.org/abs/2506.02379v1. A full rights notice is delivered at licenses/arxiv-2506.02379-CC-BY-4.0.txt. Characters: every retained excerpt is pure ASCII, so the delivered engine, XeLaTeX with its default Latin Modern text font, reports no missing character.
\begin{prop}\label{prop: coincidence of p'}
  Let $m,n \in \mathbb{Z}_{\geq 0}$, $\alpha_i,\beta_i,\gamma_j,\delta_j \in \mathbb{C}$ $(i=1,\dots,m,j=1,\dots,n)$.
  If
  \[
    p'_k(\alpha_1,\dots,\alpha_m;\beta_1,\dots,\beta_m) = p'_k(\gamma_1,\dots,\gamma_n;\delta_1,\dots,\delta_n) \ \text{ for all } k \in \mathbb{Z}_{> 0},
  \]
  then, we have
  \[
    \sum_{\substack{1 \leq i \leq m\\\alpha_i=\epsilon}} \beta_i = \sum_{\substack{1 \leq j \leq n\\\gamma_j=\epsilon}} \delta_j \ \text{ for all } \epsilon \in \mathbb{C}.
  \]
  In particular, if
  \[
    p_k(\alpha_1,\dots,\alpha_m) = p_k(\gamma_1,\dots,\gamma_n) \ \text{ for all } k \in \mathbb{Z}_{> 0},
  \]
  then we have
  \[
    \sharp\{ i \mid \alpha_i = \epsilon \} = \sharp\{ j \mid \gamma_j = \epsilon \} \ \text{ for all } \epsilon \in \mathbb{C}^\times.
  \]
\end{prop}

\begin{thm}\label{thm: sufficient condition for evaluation2}
  Suppose that $A \setminus \{0\}$ is closed under the addition.
  Let $(c_\mathbf{a})_{\mathbf{a} \in (A \setminus \{0\})^*}$ be a family of complex numbers satisfying the following for all $r \in \mathbb{Z}_{\geq 0}$ and $\mathbf{a} = (a_1,\dots,a_r) \in (A \setminus \{0\})^r$$:$
  \begin{align*}
    &c_{()} = 1,\\
    &c_\mathbf{a} = 0 \ \text{ if } r > n,\\
    &c_a c_\mathbf{a} = c_{a*\mathbf{a}} + \sum_{l=1}^r c_{(a_1,\dots,a_l+a,\dots,a_r)} \ \text{ for all } a \in A \setminus \{0\}.
  \end{align*}
  Then, there exists $\mathbf{p} \in X^n$ such that $m_\mathbf{a}(\mathbf{p}) = c_\mathbf{a}$ for all $\mathbf{a} \in (A \setminus \{0\})^*$.
\end{thm}

\begin{thm}\label{thm: classification AI1}
  For each $\phi \in (\mathcal{L}^{\mathrm{tw},0})^*$, the following are equivalent:
  \begin{enumerate}
    \item $\dim V(\phi) < \infty$.
    \item There exist $\nu_1,\nu_{-1} \in \mathbb{C}$, $n \in \mathbb{Z}_{\geq 0}$, and $\alpha_1,\dots,\alpha_n \in \mathbb{C}^\times \setminus \{ \pm 1 \}$ such that
    \[
      \phi(w_{a,0}) = \nu_1 \cdot 2^a + \nu_{-1} \cdot (-2)^a + \sum_{i=1}^n (\alpha_i+\alpha_i^{-1})^a \ \text{ for all } a \geq 0.
    \]
  \end{enumerate}
\end{thm}

\begin{lem}\label{lem: deduction from A1}
  Let $\phi \in (\mathcal{L}^{\mathrm{tw},0})^*$ be such that the highest weight module $V(\phi)$ is finite-dimensional.
  Set
  \[
    n := \max\{ r \geq 0 \mid y_{+,0,0}^{(r)} v_\phi \neq 0 \}.
  \]
  Then, there exist $\beta_1,\dots,\beta_n \in \mathbb{C}$ such that
  \[
    \phi(2w_{+,a,0}) = \sum_{i=1}^n \beta_i^a \ \text{ for all } a \in \mathbb{Z}_{\geq 0}.
  \]
\end{lem}

\begin{lem}\label{lem: deduction from AI1}
  Let $\phi \in (\mathcal{L}^{\mathrm{tw},0})^*$ be such that the highest weight module $V(\phi)$ is finite-dimensional.
  Set
  \[
    n := \max\{ r \geq 0 \mid Y_{0,1}^{(r)} v_\phi \neq 0 \}.
  \]
  Then, the following hold$:$
  \begin{enumerate}
    \item There exist $\nu_1,\nu_{-1} \in \mathbb{C}$ and $\beta_1,\dots,\beta_n \in \mathbb{C}^\times \setminus \{ \pm 2 \}$ such that
    \[
      \phi(w_{+,a,0}) = \nu_1 \cdot 2^a + \nu_{-1} \cdot (-2)^a + \sum_{i=1}^n (\beta_i+\beta_i^{-1})^a \ \text{ for all } a \in \mathbb{Z}_{\geq 0}.
    \]
    \item $\phi(d_\mathbf{a}) = 0$ for all $r > n$ and $\mathbf{a} \in B^r$.
  \end{enumerate}
\end{lem}

\begin{thm}\label{thm: classification AI2}
  For each $\phi \in (\mathcal{L}^{\mathrm{tw},0})^*$, the following are equivalent$:$
  \begin{enumerate}
    \item $\dim V(\phi) < \infty$.
    \item There exist $\nu_1,\nu_{-1} \in \frac{1}{2}\mathbb{Z}_{\geq 0}$, $n \in \mathbb{Z}_{\geq 0}$, and $\alpha_1,\dots,\alpha_r \in \mathbb{C}^\times \setminus \{ \pm 1 \}$ such that
    \begin{align*}
      &\phi(w_{+,a,0}) = \nu_1 \cdot 2^a + \nu_{-1} \cdot (-2)^a + \sum_{i=1}^n (\alpha_i+\alpha_i^{-1})^a,\\
      &\phi(w_{-,a,1}) = \sum_{i=1}^n (\alpha_i+\alpha_i^{-1})^a (\alpha_i-\alpha_i^{-1}).
    \end{align*}
  \end{enumerate}
\end{thm}

\begin{proof}
  First, assume the condition $(2)$, and consider the tensor product module
  \[
    V := V_\mathfrak{k}(\nu_1)_1 \otimes V_\mathfrak{k}(\nu_{-1})_{-1} \otimes V(1)_{\alpha_1} \otimes \cdots \otimes V(1)_{\alpha_n}.
  \]
  Then, the submodule of $V$ generated by $v_{\nu_1} \otimes v_{\nu_{-1}} \otimes v_1 \otimes \dots \otimes v_1$ is a highest weight module of highest weight $\phi$.
  This implies that $V(\phi)$ is a quotient of $V$, and hence, finite-dimensional.

  Next, assume the condition $(1)$.
  Set
  \[
    r_0 := \max\{ r \geq 0 \mid y_{+,0,0}^{(r)} v_\phi \neq 0 \}.
  \]
  By Lemma \ref{lem: deduction from A1}, there exist $\beta'_1,\dots,\beta'_{r_0} \in \mathbb{C}$ such that
  \begin{align}\label{eq: phi(2w+a0) from A1}
    \phi(2w_{+,a,0}) = \sum_{i=1}^{r_0} (\beta'_i)^a \ \text{ for all } a \in \mathbb{Z}_{\geq 0}.
  \end{align}

  Set
  \[
    r_1 := \max\{ r \geq 0 \mid Y_{0,1}^{(r)} v_\phi \neq 0 \}.
  \]
  By Lemma \ref{lem: deduction from AI1}, there exist $\nu_{\pm 1} \in \mathbb{C}$ and $\beta''_1,\dots,\beta''_{r_1} \in \mathbb{C}^\times \setminus \{ \pm 2 \}$ such that
  \begin{align}\label{eq: phi(w+a0) from AI1}
    \phi(w_{+,a,0}) = \nu_1 \cdot 2^a + \nu_{-1} \cdot (-2)^a + \sum_{i=1}^{r_1} (\beta''_i)^a \ \text{ for all } a \in \mathbb{Z}_{\geq 0}.
  \end{align}

  Equations \eqref{eq: phi(2w+a0) from A1} and \eqref{eq: phi(w+a0) from AI1} imply that
  \[
    p'_k(\beta'_1,\dots,\beta'_{r_0};\beta'_1,\dots,\beta'_{r_0}) = p'_k(2,-2,\beta''_1,\dots,\beta''_{r_1};4\nu_1,-4\nu_{-1},2\beta''_1,\dots,2\beta''_{r_1}) \ \text{ for all } k \geq 1.
  \]
  By Proposition \ref{prop: coincidence of p'}, we obtain
  \begin{align*}
    &\pm 2 \cdot \sharp\{ i \mid \beta'_i = \pm 2 \} = \pm 4\nu_{\pm 1},\\
    &\epsilon \cdot \sharp\{ i \mid \beta'_i = \epsilon \} = 2\epsilon \cdot \sharp\{ i \mid \beta''_i = \epsilon \} \ \text{ for all } \epsilon \in \mathbb{C}^\times \setminus \{ \pm 2 \}.
  \end{align*}
  In particular,
  \[
    \nu_{\pm 1} \in \frac{1}{2}\mathbb{Z}_{\geq 0}.
  \]

  By Lemma \ref{lem: deduction from AI1}, we have
  \begin{align}\label{eq: phi(dbfa) = 0 AI2}
    \phi(d_\mathbf{a}) = 0 \ \text{ for all } r > r_1,\ \mathbf{a} \in B^r.
  \end{align}
  We will show that
  \[
    \phi(d_\mathbf{a}) = 0 \ \text{ for all } r > r_1,\ \mathbf{a} \in (\mathbb{Z}_{\geq 0} \times \mathbb{Z}_{> 0})^r.
  \]
  Let $r > r_1$ and $s \geq 0$ be such that $\phi(d_\mathbf{a}) = 0$ for all $\mathbf{a} = ((a_1,b_1),\dots,(a_r,b_r))$ with $\sharp\{i \mid \text{$b_i$ is even }\} \leq s$.
  Then, for each $(a,b) \in \mathbb{Z}_{\geq 0} \times \mathbb{Z}_{> 0,\mathrm{even}}$, we have
  \[
    \phi(d_{(a,b)*\mathbf{a}}) = \phi(d_{(a,b)}) \phi(d_\mathbf{a}) - \sum_{l=1}^r \phi(d_{((a_1,b_1),\dots,(a_l+a,b_l+b),\dots,(a_r,b_r))}) = 0.
  \]
  Hence, $\phi(d_\mathbf{b}) = 0$ for all $\mathbf{b} = ((a_1,b_1),\dots,(a_{r+1},b_{r+1}))$ with $\sharp\{i \mid \text{$b_i$ is even }\} \leq s+1$.
  Also, let $q \geq 0$ be such that $\phi(d_\mathbf{a}) = 0$ for all $\mathbf{a} = ((a_1,b_1),\dots,(a_r,b_r))$ with
  \[
    \sharp\{i \mid \text{$b_i$ is even }\} \leq s+1 \text{ and } \sharp\{ i \mid (a_i,b_i) = (0,1) \} = r-q > 0.
  \]
  We may assume that $(a_1,b_1),\dots,(a_q,b_q) \neq (0,1)$ and $(a_{q+1},b_{q+1}),\dots,(a_r,b_r) = (0,1)$.
  Then, for each $(a,b) \in \mathbb{Z}_{\geq 0} \times \mathbb{Z}_{> 0,\mathrm{odd}}$, we have
  \begin{align*}
    0 = \phi(d_{(a,b)*\mathbf{a}}) &= \phi(d_{(a,b)}) \phi(d_\mathbf{a}) - \sum_{l=1}^r \phi(d_{((a_1,b_1),\dots,(a_l+a,b_l+b),\dots,(a_r,b_r))})\\
    &= -(r-q)\phi(d_{((a_1,b_1),\dots,(a_q,b_q),(a,b+1),(0,1),\dots,(0,1))}).
  \end{align*}
  This implies that $\phi(d_\mathbf{a}) = 0$ for all $\mathbf{a} = ((a_1,b_1),\dots,(a_r,b_r))$ with
  \[
    \sharp\{i \mid \text{$b_i$ is even }\} \leq s+1 \text{ and } \sharp\{ i \mid (a_i,b_i) = (0,1) \} = r-q-1.
  \]
  By above, our claim follows by induction.

  For each $r \in \mathbb{Z}_{\geq 0}$ and $\mathbf{a} = ((a_1,b_1),\dots,(a_r,b_r)) \in (A \setminus \{(0,0)\})^r$, define $c_\mathbf{a}$ by
  \[
    c_\mathbf{a} := \phi(d_{((a_1,a_1+b_1),\dots,(a_r,a_r+b_r))}).
  \]
  By the definition of $d_{\mathbf{a}}$'s, equation \eqref{eq: phi(dbfa) = 0 AI2}, and Theorem \ref{thm: sufficient condition for evaluation2}, there exist $\delta_1,\dots,\delta_{r_1},\gamma_1,\dots,\gamma_{r_1} \in \mathbb{C}$ such that
  \[
    c_\mathbf{a} = m_\mathbf{a}((\delta_1,\gamma_1),\dots,(\delta_{r_1},\gamma_{r_1})) \ \text{ for all } \mathbf{a} \in (A \setminus \{(0,0)\})^*.
  \]
  As in the proof of Theorem \ref{thm: classification AI1}, we obtain
  \[
    \delta_i^2 = (\gamma_i^2+4)\gamma_i^2 \ \text{ for all } i = 1,\dots,r_1.
  \]

  For each $a \geq 1$, we have
  \[
    \phi(w_{+,0,2a}) = \phi(d_{(0,2a)}) = c_{(0,2a)} = \sum_{i=1}^r \gamma_i^{2a}.
  \]
  On the other hand, equation \eqref{eq: phi(w+a0) from AI1} implies that
  \[
    \phi(w_{+,0,2a}) = \sum_{i=1}^{r_1} ((\beta''_i)^2-4)^a \ \text{ for all } a \geq 1.
  \]
  By reordering the tuple $((\delta_1,\gamma_1),\dots,(\delta_{r_1},\gamma_{r_1}))$, we may assume that
  \[
    \gamma_i^2 = (\beta''_i)^2-4 \neq 0 \ \text{ for all } i = 1,\dots,r_1.
  \]

  Set $\beta_i := \delta_i \gamma_i^{-1}$.
  Then, we obtain
  \begin{align}\label{eq: beta2 = gamma2+4 neq 4}
    \beta_i^2 = \gamma_i^2+4 = (\beta''_i)^2 \neq 4.
  \end{align}

  As in the proof of Theorem \ref{thm: classification AI1}, for each $a \in \mathbb{Z}_{\geq 0}$, we have
  \begin{align*}
    \phi(w_{+,2a,0}) &= \sum_{i=1}^{r_1} \beta_i^{2a} + 4^a(\phi(w_{+,0,0})-r_1) \\
    &= \sum_{i=1}^{r_1} \beta_i^{2a} + 4^a(\nu_1+\nu_{-1}) \\
    \phi(w_{+,2a+1,0}) &= \sum_{i=1}^{r_1} \beta_i^{2a+1} + 4^a(\phi(w_{+,1,0})-\sum_{i=1}^{r_1} \beta_i).
  \end{align*}
  On the other hand, equation \eqref{eq: phi(w+a0) from AI1} implies that
  \begin{align*}
    \phi(w_{+,2a+1,0}) = 2 \cdot 4^a(\nu_1 - \nu_{-1}) + \sum_{i=1}^{r_1} (\beta''_i)^{2a+1} = 2 \cdot 4^a(\nu_1 - \nu_{-1}) + \sum_{i=1}^{r_1} \beta''_i \beta_i^{2a}.
  \end{align*}
  Comparing this and the previous identity, we obtain $\sum_{i=1}^{r_1} (\beta_i-\beta''_i) = 0$.
  Therefore,
  \[
    \phi(w_{+,2a+1,0}) = \sum_{i=1}^{r_1} \beta_i^{2a+1} + 2 \cdot 4^a(\nu_1-\nu_{-1}).
  \]
  
  Similarly, we have
  \begin{align*}
    &\phi(w_{-,2a,1}) = -\sum_{i=1}^{r_1} \beta_i^{2a}\gamma_i,\\
    &\phi(w_{-,2a+1,1}) = -\sum_{i=1}^{r_1} \beta_i^{2a+1}\gamma_i,
  \end{align*}
  for all $a \in \mathbb{Z}_{\geq 0}$.
  Summarizing, we have obtained
  \begin{align*}
    &\phi(w_{+,a,0}) = \nu_1 \cdot 2^a + \nu_{-1} \cdot (-2)^a + \sum_{i=1}^{r_1} \beta_i^a,\\
    &\phi(w_{-,a,1}) = -\sum_{i=1}^{r_1} \beta_i^a \gamma_i,
  \end{align*}
  for all $a \in \mathbb{Z}_{\geq 0}$.
  For each $i = 1,\dots,r_1$, there exists $\alpha_i \in \mathbb{C}^\times$ such that
  \[
    \beta_i = \alpha_i + \alpha_i^{-1} \text{ and } \gamma_i = -(\alpha_i-\alpha_i^{-1}).
  \]
  The condition \eqref{eq: beta2 = gamma2+4 neq 4} implies that $\alpha_i \neq \pm 1$.
  Thus, we complete the proof.
\end{proof}

\end{document}
